\documentclass[a4paper]{umbruch-position}
\usepackage[utf8]{inputenc}
\usepackage[T1]{fontenc}
\usepackage[ngerman]{babel}
\usepackage{tikz}
\usepackage{amssymb}
\usepackage{eurosym}
\usepackage{graphicx}
\usepackage{xcolor}
\definecolor{umbruch-dark}{HTML}{2A2A2A}
\definecolor{umbruch-teal}{HTML}{008080}
\definecolor{umbruch-red}{HTML}{B22222}
\usetikzlibrary{shapes.geometric, arrows.meta, positioning, fit, calc}

\positionsnummer{Basiert auf der Studie ``Systemtheoretische Aufarbeitung der Regressangst'' (ST-001)}
\versionsnummer{0.2}
\title{Die Regress-Firewall}
\untertitel{Was niedergelassene \"Arzte wissen m\"ussen --- und was sie tun k\"onnen,\\bevor es zum Verfahren kommt}

\begin{document}

\maketitle

\begin{infobox}
\textbf{WICHTIGER HINWEIS:}\\
\textit{Dieser Leitfaden dient der allgemeinen Information und stellt keine Rechtsberatung dar. F\"ur individuelle rechtliche Beratung wenden Sie sich an einen Fachanwalt f\"ur Medizin- oder Sozialrecht oder Ihre jeweilige Kassen\"arztliche Vereinigung (KV).}
\end{infobox}

\section{Die Zahlen zur Regress-Bedrohung}

Viele niedergelassene \"Arzte f\"urchten Wirtschaftlichkeitspr\"ufungen so sehr, dass sie im Zweifel notwendige Verordnungen meiden (Quartalsende-Effekt). Doch betrachten wir die statistische Realit\"at:

\begin{itemize}
    \item \textbf{Regressquote:} Die Wahrscheinlichkeit eines durchgesetzten Regresses liegt bei unter \textbf{1\,\%} (V1+V2 zusammen). Beispiel KV BaW\"u 2022: 21 von ca. 21.500 \"Arzten = 0,097\,\%.
    \item \textbf{Extremfall:} BSG 2025 (B\,6\,KA\,9/24\,R): 491.000\,EUR Regress wegen Stempel statt Unterschrift --- bei medizinisch korrekten Verordnungen.
    \item \textbf{Gef\"uhlte Bedrohung:} 72\,\% der Haus\"arzte berichten eigene Regresserfahrung (Ribbat 2023, n=770, V1+V2 zusammen). 46,7\,\% der Medizinstudierenden nennen Regresse als Niederlassungshindernis (KBV 2018).
\end{itemize}

Die \"offentlich kommunizierte 1-\%-Zahl vermischt vier grundlegend verschiedene Messebenen:

\begin{center}
{\small
\begin{tabular}{p{3.5cm} p{7cm}}
\toprule
\textbf{Messebene} & \textbf{Was sie misst} \\
\midrule
Pr\"ufquote & Wieviele Auff\"alligkeiten entstehen (Screening-Output) \\
Verfahrensquote & Wieviele Einzelfallpr\"ufungen eingeleitet werden \\
Durchsetzungsquote & Wieviele Regresse tats\"achlich festgesetzt werden \\
Volumenquote & Wieviel Geld tats\"achlich eingezogen wird \\
\bottomrule
\end{tabular}
}
\end{center}

\begin{center}
  \includegraphics[width=0.90\textwidth]{C:/Users/User/OneDrive/.UMBRUCH/Projekte/Regress-Melder/grafiken/vier_messebenen}
\end{center}

\textbf{Die 1-\%-Zahl gilt f\"ur das Screening.} F\"ur Einzelfallverfahren gibt es keine bundesweit ver\"offentlichten Zahlen~--- und die Erfolgsquote der Kassen ist dort erheblich h\"oher. Pr\"avention setzt deshalb nicht beim Screening an, sondern bei der Einzelfall-Angreifbarkeit.

Das Vermeiden einer leitliniengerechten Therapie aufgrund unbegr\"undeter Budget-\"Angste schadet Patienten. \textbf{Es gibt Wege,} sich rechtssicher vorab im System zu bewegen.

\section{H\"aufige Ausl\"oser f\"ur Regresse (§ 106 SGB V)}

Gem\"a\ss{} § 12 SGB V m\"ussen die Leistungen ausreichend, zweckm\"a\ss{}ig und wirtschaftlich sein.

\begin{infobox}
\textbf{1. Auff\"alligkeitspr\"ufung (\S\,106 SGB\,V)}\\
Ein statistischer Vergleich des Verordnungsverhaltens mit der Fachgruppe. Seit Abschaffung der bundeseinheitlichen Richtgr\"o\ss{}en (2017) gelten \textbf{regionale Pr\"ufvereinbarungen} mit eigenen Schwellenwerten (Durchschnittswerte, Zielwerte, DDD-Quoten je KV-Bezirk). Eine Auff\"alligkeitspr\"ufung droht bei signifikanter \"Uberschreitung des regionalen Schwellenwerts.
\end{infobox}

\begin{infobox}
\textbf{2. Einzelfallpr\"ufung}\\
Die Krankenkasse beantragt aus bestimmten Gr\"unden (z.\,B. Off-Label-Use) die Pr\"ufung einer ganz spezifischen Verordnung.
\end{infobox}

\begin{infobox}
\textbf{3. Heilmittel- und Formfehler-Regress}\\
Physiotherapie, Logo, Ergo: \"Anderung der Ausf\"uhrung au\ss{}erhalb des Heilmittelkatalogs oder unzul\"assige Delegation ärztlicher Leistungen. Formale Gr\"unde sind sehr h\"aufig (ohne dass die Verordnung medizinisch falsch war).
\end{infobox}

\section{Entscheidungsbaum: \glqq{}Soll ich das verordnen?\grqq{}}

\begin{infobox}
\textit{Entscheidungshilfe, keine Rechtsberatung. Fristen und Rechtsgrundlagen k\"onnen sich \"andern.}
\end{infobox}

\begin{center}
\tikzset{
    action/.style={rectangle, draw=umbruch-dark, fill=umbruch-teal!15, rounded corners, minimum height=1cm, align=center, text width=4cm, inner sep=0.2cm},
    decision/.style={diamond, aspect=2.0, draw=umbruch-dark, fill=white, align=center, text width=3cm, inner sep=0cm},
    warning/.style={rectangle, draw=umbruch-red, fill=umbruch-red!15, rounded corners, minimum height=1cm, align=center, text width=4.5cm, inner sep=0.2cm},
    arrow/.style={-{Stealth[scale=1.2]}, thick, umbruch-dark}
}
\begin{tikzpicture}[
    scale=0.8, transform shape,
    node distance=1.5cm and 1cm,
    font=\sffamily\small
]
\node[decision] (leitlinie) {Ist es leitlinien-gerecht?\\(§ 12 SGB V)};
\node[action, right=1.5cm of leitlinie] (nein_leit) {NICHT verordnen!};
\node[decision, below=1.5cm of leitlinie] (richtgrosse) {\"Uber regionalem\\Schwellenwert?\\(\S\,106 SGB V)};
\node[action, left=1.5cm of richtgrosse] (ja_verordnen1) {Verordnen,\\alles OK.};
\node[decision, below=1.5cm of richtgrosse] (besonderheit) {Praxisbesonder-heit?\\(\S\,106 SGB V)};
\node[action, right=1.5cm of besonderheit] (dok_verordnen) {VERORDNEN +\\(Detailliert in Akte)};
\node[decision, below=1.5cm of besonderheit] (kv_beratung) {KV-Beratung?\\(\S\,106b Abs. 2, nur V1)};
\node[action, left=1.5cm of kv_beratung] (kv_ja) {KV best\"atigt:\\VERORDNEN!};
\node[warning, right=1.5cm of kv_beratung] (kv_nein) {Nur bei unklarer\\Erstattungsf\"ahigkeit:\\KK-Anfrage (BSG 27/12 R)\\oder Patienteneskalation};
\draw[arrow] (leitlinie) -- node[above] {NEIN} (nein_leit);
\draw[arrow] (leitlinie) -- node[left] {JA} (richtgrosse);
\draw[arrow] (richtgrosse) -- node[above] {NEIN} (ja_verordnen1);
\draw[arrow] (richtgrosse) -- node[left] {JA} (besonderheit);
\draw[arrow] (besonderheit) -- node[above] {JA} (dok_verordnen);
\draw[arrow] (besonderheit) -- node[left] {NEIN} (kv_beratung);
\draw[arrow] (kv_beratung) -- node[above] {JA} (kv_ja);
\draw[arrow] (kv_beratung) -- node[above] {NEIN} (kv_nein);
\end{tikzpicture}
\end{center}

\newpage
\section{Die 3 wirksamsten Schutzma\ss{}nahmen (juristisch belegt)}

\begin{enumerate}
    \item \textbf{L\"uckenlose Dokumentation ZUM ZEITPUNKT der Verordnung.} SG Marburg 14.02.2024 (S 18 KA 96/23): Wer die Behandlungsdokumentation vorlegen kann, gewinnt. \textbf{Wichtig:} BSG B 6 KA 26/13 stellt klar: Nachgereichte Dokumentation reicht NICHT, sie muss zeitnah erstellt sein. Sonst greift die Beweislastumkehr nach \S\,630h BGB.
    \item \textbf{Substantiierter Vortrag von Praxisbesonderheiten BEREITS im Verwaltungsverfahren} --- nicht erst vor dem Sozialgericht (LSG BW 15.11.2023, SG Marburg 31.01.2024). KV-Kennziffer auf dem Behandlungsschein im Quartal markieren!
    \item \textbf{Vorab-Kl\"arung in den expliziten Ausnahmef\"allen:} Cannabis (\S\,31 Abs.\,6 SGB\,V), h\"ausliche Krankenpflege (\S\,37), Soziotherapie (\S\,37a), station\"are Reha (\S\,40), bestimmte Heil-/Hilfsmittel. \textbf{Nicht jedoch} bei normalen Arzneimitteln mit Indikation --- dort verbietet \S\,29 BMV-\"A die Vorabgenehmigung. In Grauzonen kann der Arzt nach BSG B 6 KA 27/12 R eine Stellungnahme der Kasse erbitten; sie bindet juristisch nicht, schafft aber Beweislast-Vorteile.
\end{enumerate}

\section{Vier Schutzinstrumente: Zwei werden fast nie genutzt}

Das System kennt vier Ausgleichsmechanismen~--- aber nur zwei davon sind allgemein bekannt:

\begin{center}
{\small
\begin{tabular}{p{2.5cm} p{3cm} p{5.5cm}}
\toprule
\textbf{Wann?} & \textbf{Instrument} & \textbf{Was es bringt} \\
\midrule
Vor Verordnung & \textbf{KV-Verordnungs\-beratung} & Fachlich gest\"utzte Ent\-scheidungs\-grundlage; st\"arkt Ihre Dokumentation erheblich \\[4pt]
Vor Verordnung & \textbf{KK-Anfrage zur Verordnungs\-f\"ahigkeit} (BSG B\,6\,KA\,27/12\,R) & Arzt fragt Kasse um Stellungnahme (keine Genehmigung, daher nicht durch \S\,29 BMV-\"A blockiert). \textbf{Sinnvoll bei unklarer Erstattungsf\"ahigkeit} (Off-Label, neue Wirkstoffe). Bei Standardmedikamenten inhaltsleer (Verordnungsf\"ahigkeit steht ja fest). Positive Antwort: Schutz vor V2-Regress \\[4pt]
Nach Auff\"alligkeit & Praxisbesonder\-heiten & Statistisches Signal wird korrigiert --- aber nur wenn aktiv geltend gemacht \\[4pt]
Nach Bescheid & Beratung vor Regress (\S\,106b Abs.\,2) & P\"adagogischer Schritt statt Sofort\-regress --- gilt nur bei Auff\"alligkeitspr\"ufung, nicht Einzelfall \\
\bottomrule
\end{tabular}
}
\end{center}

\begin{infobox}
\textbf{Wichtige Einschr\"ankungen KK-Anfrage (BSG 27/12\,R):}\\
(1)~Positive Antwort sch\"utzt vor KK-initiiertem Einzelfallregress (V2) --- \textbf{nicht} vor Auff\"alligkeitspr\"ufung (V1).\\
(2)~Die Anfrage ist eine \textit{Stellungnahme}, keine Genehmigung --- daher \textbf{nicht} durch \S\,29 BMV-\"A blockiert (BSG Rn.\,19).\\
(3)~Bei \textbf{klar verordnungsf\"ahigen Standardmedikamenten} ist die Anfrage inhaltsleer: Die Verordnungsf\"ahigkeit steht fest, die Kasse hat keinen Anlass zu antworten.\\
(4)~\textbf{Sinnvoll} bei: Off-Label-Use, neuen Wirkstoffen, grundsätzlich ausgeschlossenen Arzneimitteln, unklarer Indikation.
\end{infobox}

\section{H\"aufige Fehler vermeiden}

\begin{enumerate}
    \item \textbf{Am Quartalsende sparen:} Teure, notwendige Behandlungen zu verschieben bringt juristisch keinen Vorteil, sch\"adigt aber den Patienten.
    \item \textbf{Praxisbesonderheiten nicht geltend machen:} Dokumentieren Sie exzessiv teure Zyklen (z.\,B. hohe Facharztdichte am Land, hochselektioniertes Patientenkollektiv).
    \item \textbf{Falsche Dokumentation bei Delegation:} Oft fallen \"Arzte beim Heilmittelregress in die Falle der unzul\"assigen Delegation (Formfehler!), weil sie die Vorgaben formal falsch per EBM kodieren.
    \item \textbf{\"Uberweisen statt verordnen:} Viele Haus\"arzte \"uberweisen an Spezialisten, was eine \"Uberweisungskaskade und unn\"otige Kosten nach sich zieht.
\end{enumerate}

\section{Die zwei Pr\"ufdom\"anen: Wo das echte Risiko liegt}

\textbf{Wichtig vorab:} Das Pr\"ufsystem misst prim\"ar formale und dokumentarische Korrektheit, nicht medizinische Qualit\"at. Allerdings: V1 ber\"ucksichtigt medizinische Besonderheiten im Schutzstufenmodell (Anh\"orung, Praxisbesonderheiten). V2a (unwirtschaftliche Verordnung) l\"asst medizinische Verteidigung zu. Nur V2b (Formfehler) ist medizinisch irrelevant (doppelter Zertifikatsfehler, ST-001 Abschn.~2.5). \textbf{Zugleich gilt:} Gesch\"atzt ${\sim}$35\,\% der V2-Pr\"ufungen codieren medizinische Evidenz als Formregeln (Biosimilar-Pflicht, AM-RL, G-BA-Beschl\"usse) --- sie sind keine Willk\"ur, sondern Evidenzsteuerung. Das Problem ist nicht die Steuerung selbst, sondern dass sie punitiv statt edukativ erfolgt: Der Arzt erh\"alt kein Lernsignal, nur ein Strafsignal. \textbf{Ihre Pr\"avention muss daher auf Dokumentation und Formkorrektheit zielen --- nicht auf medizinische Zur\"uckhaltung.}

\textbf{Positive Merkmale der Auff\"alligkeitspr\"ufung (V1):} Das V1-Verfahren enth\"alt ein Schutzstufenmodell mit mehreren Entlastungsstufen. Erstmalig auff\"allige \"Arzte erhalten eine \textbf{Karenzzeit nach Erstberatung} (n\"achster Pr\"ufzeitraum bleibt ungepr\"uft). In den ersten zwei Jahren nach Beratung gilt ein \textbf{Regressdeckel von 25.000\,\euro{}}. Zudem greift die \textbf{5-Jahres-Regel:} Wer f\"unf Jahre lang unauff\"allig bleibt, gilt wieder als erstmalig auff\"allig --- mit erneutem Anspruch auf Beratung vor Regress.

Pr\"avention muss zwischen zwei grundlegend verschiedenen Verfahren unterscheiden:

{\small
\begin{tabularx}{\textwidth}{p{3cm} X X}
\toprule
\textbf{Kategorie} & \textbf{Auff\"alligkeitspr\"ufung} & \textbf{Einzelfallpr\"ufung} \\
\midrule
\multicolumn{3}{l}{\textit{Verfahrensstruktur}} \\
\addlinespace[2pt]
Ausl\"oser & Abweichung vom Fachgruppen\-durch\-schnitt & Konkreter Kassenantrag \\
\addlinespace
Rechts\-grundlage & \S\,106 SGB\,V & \S\,106b Abs.\,1 SGB\,V \\
\addlinespace
Charakter & Such- und Filterverfahren & Gezieltes Verdachts\-verfahren \\
\addlinespace
Verdacht & Kein (statistisches Signal) & Ja (individueller Vorwurf) \\
\midrule
\multicolumn{3}{l}{\textit{Schutzstruktur}} \\
\addlinespace[2pt]
Beratung vor Regress & Vorgeschrieben & Gesetzlich ausgeschlossen \\
\addlinespace
Pr\"avention & Praxis\-besonderheiten, Dokumentation & Option\,B (KV-Beratung), Option\,C (KK-Anfrage), antragsfeste Doku\-mentation \\
\midrule
\multicolumn{3}{l}{\textit{Risikoprofil}} \\
\addlinespace[2pt]
Betroffenheit & Teilweise & 100\,\% \\
\addlinespace
Neg.\,Ausgang & Niedrig wahrscheinlich & Hoch wahrscheinlich \\
\addlinespace
Schaden\-h\"ohe & Niedrig bis moderat & Moderat bis hoch \\
\addlinespace
Wiederholbarkeit & Eingeschr\"ankt & M\"oglich \\
\addlinespace
Existenz\-bedrohend & Sehr gering & Niedrig, aber real \\
\bottomrule
\end{tabularx}
}

\textbf{Konsequenz:} Pr\"avention setzt nicht bei Auff\"alligkeiten an. Sie setzt bei der formalen und dokumentarischen Angreifbarkeit an~--- denn das ist, was im Einzelfallverfahren gepr\"uft wird.

\medskip
\textbf{Wo gilt die 1-\,\%-Zahl?} Beide Verfahren als Eskalationstreppe~--- von unten lesen:

{\footnotesize
\begin{tabularx}{\textwidth}{p{3.0cm} p{0.5cm} p{1.8cm} p{1.8cm}}
\toprule
\textbf{Schritt} & \textbf{R.} & \textbf{V1: Auffälligkeit} & \textbf{V2: Einzelfall} \\
\midrule
Schritt n: Regress rechts\-kr\"aftig & 1 & k.D.* (KV BaW\"u: 0,097\,\%, V1+V2) & k.D.* \\
\addlinespace
Schritt 5: Regress fest\-gesetzt & 1 & $\sim$0,1\,\% & k.D. --- hoch \\
\addlinespace
Schritt 4: Entscheidung & 0 & hoch ($>$99\,\% kein Regress) (Ergebnis des Schutzstufenmodells: Anh\"orung, Praxisbesonderheiten, Beratung) & mittel--hoch \\
\addlinespace
Schritt 3: Beratung vor Regress & 0 & vorgeschrieben & \textit{fehlt gesetzlich} \\
\addlinespace
Schritt 2: Anh\"orung / Pr\"ufung & 0 & $\sim$100\,\% & $\sim$100\,\% \\
\addlinespace
Schritt 1: Pr\"uf\-verfahren ein\-geleitet & 0 & $\sim$10--20\,\% [Sch.] & k.D. \\
\addlinespace
Schritt 0: Trigger & 0 & Abweichung erkannt & Kassen\-antrag \\
\bottomrule
\end{tabularx}

\smallskip
\noindent R.\,=\,Regress eingetreten (0/1). \quad k.D.\,=\,keine publizierten Daten. \quad Sch.\,=\,Sch\"atzung.\\
{[*] Quellen: (1)~KV\,BaW\"u 2022: 0,097\,\% (21\,/\,21.500\,\"Arzte, V1+V2; kvbawue.de). (2)~KV\,Hessen 2015: 3\,/\,916 = 0,3\,\%\,$|$\,S0 (\textit{Dt.~\"Arzteblatt} Art.~113630; V1). Die ,,unter 1\,\%``-Zahl gilt V1+V2 zusammen~--- keine Verfahrenstrennung ver\"offentlicht.}
}

\begin{center}
  \includegraphics[width=\textwidth]{C:/Users/User/OneDrive/.UMBRUCH/Projekte/Regress-Melder/grafiken/pipeline_vergleich_v1_v2}
  \smallskip\\
  {\small\textit{V1 = Schwellenwert-Screening (viele markiert, fast alle durch Schutzstufenmodell entlastet). V2b = nahezu deterministisch bei Formfehlern. V2a = Verteidigung m\"oglich.}}
\end{center}

\section{Wichtig: Auch unauff\"alliger Gesamtbefund sch\"utzt nicht}

Die statistische Auff\"alligkeitspr\"ufung (\S\,106 SGB\,V) und die Einzelfallpr\"ufung (\S\,106b SGB\,V) sind zwei \textit{unabh\"angige} Verfahren. Auch bei einem unauff\"alligen statistischen Gesamtbild kann eine Krankenkasse eine Einzelfallpr\"ufung beantragen. Das ist die wichtigste Praxis-Erkenntnis: Statistische Wirtschaftlichkeit sch\"utzt nicht vor Einzelfallregress. Allein in der KV Nordrhein werden etwa 20.000 Einzelfallpr\"ufungs-Antr\"age pro Jahr gestellt --- bei einem statistisch ``unauff\"alligen'' Korpus von 13.500 Praxen.

\section{F\"alle aus dem Leben (Anonymisiert)}

\textbf{Fall A: Physiotherapie-Regress (65.000 EUR, Dortmund):}
Ein Allgemeinarzt erhielt eine R\"uckforderung \"uber 65.000 EUR. Regressausl\"oser: Die Heilmittelkosten waren doppelt so hoch wie der Fachgruppendurchschnitt. Er vers\"aumte es, fr\"uhzeitig die \glqq{}Praxisbesonderheiten\grqq{} (viele gel\"ahmte Patienten) transparent zu codieren. 

\textbf{Fall B: Stempel-Regress --- 491.000 EUR (BSG B\,6\,KA\,9/24\,R, 27.08.2025):}
Ein niedergelassener Arzt verwendete auf seinen Rezepten einen Namensstempel statt der pers\"onlichen Unterschrift. Alle Verordnungen waren medizinisch korrekt und leitliniengerecht. Das BSG best\"atigte den vollen Regress \"uber 491.000\,EUR --- das Argument der medizinischen Korrektheit wurde als \textit{irrelevant} verworfen. Der normative Schadensbegriff greift: Der Schaden der Kasse gilt ab dem Moment der Bezahlung der formal fehlerhaften Verordnung. \textbf{Der fatale Fehler:} Ein Formcheck-Prozess im Backoffice (Vier-Augen-Prinzip) h\"atte den Stempel bei jeder einzelnen Verordnung erkannt.

\newpage
\section{Die zentrale Pr\"aventionsstrategie: Formfehler auf Null setzen}

\textbf{F\"ur V1:} Praxisbesonderheiten aktiv und fr\"uhzeitig anmelden. KV-Beratung nutzen. Das Schutzstufenmodell greift --- aber nur wenn Sie es kennen und nutzen.

\textbf{F\"ur V2b:} Das Pr\"ufsystem misst nicht medizinische Qualit\"at, sondern formale Korrektheit (doppelter Zertifikatsfehler, siehe Abschnitt~6). Ein Formfehler, der einmal im System ist, kann nach Verfahrenseinleitung nicht mehr geheilt werden (normativer Schadensbegriff, BSG B\,6\,KA\,5/09\,R). Die logische Konsequenz:

\begin{infobox}
\textbf{Kernbotschaft:} Bauen Sie eine \textbf{Firewall gegen Einzelfallpr\"ufungen} auf. Machen Sie gute Medizin --- und lassen Sie Ihre Organisation daf\"ur sorgen, dass die Dokumentation wasserdicht ist. \textbf{Kein Formfehler darf das Haus verlassen.}
\end{infobox}

\subsection{Pr\"avention als Organisationsaufgabe}

Regress-Pr\"avention ist \textbf{keine \"arztliche Aufgabe, sondern eine organisatorische}. Der Arzt soll verordnen, was medizinisch richtig ist. Die Praxisorganisation soll sicherstellen, dass die Verordnung formal korrekt codiert, vollst\"andig ausgef\"ullt und unterschrieben ist.

\textbf{Konkret:}
\begin{enumerate}[nosep]
  \item \textbf{Geschultes Backoffice-Personal:} MFA oder Praxismanager, die Verordnungen vor Ausgabe auf Vollst\"andigkeit pr\"ufen (Unterschrift, Datum, LANR/BSNR, Dosierungsangabe, ICD-Code).
  \item \textbf{Checklisten-Prozess:} Jede Verordnung durchl\"auft einen kurzen Formcheck bevor sie die Praxis verl\"asst --- analog zur Vier-Augen-Pr\"ufung in der Luftfahrt.
  \item \textbf{PVS-Autokorrektur:} Praxisverwaltungssoftware so konfigurieren, dass Pflichtfelder nicht leer bleiben k\"onnen. Viele PVS bieten Warnmeldungen bei fehlenden Angaben --- diese \textbf{nicht abschalten}.
  \item \textbf{Quartalsselbstaudit:} Stichprobenartig 20 Verordnungen des letzten Quartals auf Formfehler pr\"ufen. Fehler, die vor Pr\"ufverfahren entdeckt werden, sind heilbar.
  \item \textbf{Fallback-Prinzip:} Kein System ist perfekt. Deshalb: Dokumentation der medizinischen Begr\"undung in der Patientenakte --- als zus\"atzliche Verteidigungslinie f\"ur den Fall, dass ein Formfehler durchrutscht.
\end{enumerate}

\textbf{Der Effekt:} Der Arzt wird entlastet und kann sich auf das konzentrieren, wof\"ur er ausgebildet wurde --- gute Medizin. Die Organisation \"ubernimmt die formale Absicherung. Das ist moralisch nicht verwerflich, sondern die rationale Antwort auf ein System, das nur formale Korrektheit pr\"uft.

\subsection{Gut zu wissen: Wie weit zur\"uck d\"urfen die Sie pr\"ufen?}

Seit der Reform durch TSVG (2019) und GVWG (2021) gelten verk\"urzte \textbf{Ausschlussfristen} (das Recht erlischt automatisch --- Sie m\"ussen nichts einwenden):

\begin{itemize}[nosep]
  \item \textbf{V1 (Auff\"alligkeitspr\"ufung):} \textbf{2 Jahre} ab Jahresende (\S\,106 Abs.\,3 S.\,3 SGB\,V)
  \item \textbf{V2 (Einzelfallpr\"ufung):} \textbf{2,5 Jahre} (18 Mon. Antragsfrist + 12 Mon. Festsetzung, \S\,106 Abs.\,3 S.\,4--5 SGB\,V)
  \item \textbf{Altf\"alle (vor 11.05.2019):} \textbf{4 Jahre} (BSG 6\,RKa\,40/94, 20.09.1995)
\end{itemize}

\textbf{Achtung:} Wurde ein Verfahren rechtzeitig eingeleitet, gibt es f\"ur den anschlie\ss{}enden Instanzenzug \textbf{keine Obergrenze}. Verfahren mit 6--11 Jahren Gesamtdauer sind dokumentiert. Die Pr\"avention (kein Formfehler, kein Verfahren) bleibt der einzig sichere Schutz.

\section{Checkliste: Quartal sicher abschlie\ss{}en}

\begin{center}
\textbf{Bitte zur Erinnerung ins Backoffice h\"angen!}
\end{center}

\begin{itemize}
    \item[$\square$] \textbf{Formcheck-Prozess aktiv?} (Wird jede Verordnung vor Ausgabe auf Vollst\"andigkeit gepr\"uft?)
    \item[$\square$] \textbf{Alle Verordnungen pers\"onlich unterschrieben?} (Kein Stempel! BSG 2025: 491.000\,EUR Regress.)
    \item[$\square$] \textbf{Praxisbesonderheiten gemeldet?} (Ist mein Patientenklientel besonders morbide? Sind teure Langzeittherapien exakt codiert?)
    \item[$\square$] \textbf{Hochpreisige Verordnungen dokumentiert?} (Medizinische Begr\"undung in der Akte!)
    \item[$\square$] \textbf{Delegationen korrekt ausgef\"ullt?} (Kreuzchen, Stempel \textbf{und} Unterschrift!)
    \item[$\square$] \textbf{PVS-Warnmeldungen aktiv?} (Pflichtfelder nicht abgeschaltet?)
    \item[$\square$] \textbf{Stichproben-Selbstaudit durchgef\"uhrt?} (20 Verordnungen gepr\"uft?)
    \item[$\square$] \textbf{Beratungsgespr\"ach wahrgenommen?} (\S\,106b Abs.\,2 SGB\,V --- nur bei Auff\"alligkeitspr\"ufung!)
\end{itemize}

\vspace{1cm}

\begin{infobox}
\textbf{Gemeinsamer Rahmen (gilt f\"ur alle drei Leitf\"aden):}\\
Die ``unter 1\,\%''-Regressquote z\"ahlt nur V1-Regresse, nur rechtskr\"aftig = Schicht~1. Bundesweit gab es 2022 aber ca.\,47.000 V2-Verfahren (BMG) = Schicht~3. Der Faktor: ca.\,500$\times$. Beide Zahlen stimmen --- sie messen verschiedene Dinge (Definitionsdivergenz). In Bayern: 13.332 festgesetzte Regresse bei $\sim$24.000 \"Arzten (SpiFa 2024), Durchschnitt 334\,EUR. \textbf{Kleines Risiko pro Fall summiert sich:} Ein Arzt mit 3.000 Verordnungen pro Quartal hat \"uber 10 Jahre 120.000 Angriffspunkte. Das erkl\"art Ribbat: 72\,\% der Haus\"arzte berichten eigene Regresserfahrung. Die Angst ist rational. Sicherheit entsteht durch Formfehler-Pr\"avention --- nicht durch medizinische Zur\"uckhaltung.
\end{infobox}

\vfill
\begin{center}
\footnotesize{\textit{Dieser Leitfaden dient der allgemeinen Information und stellt keine Rechtsberatung dar. F\"ur individuelle rechtliche Beratung wenden Sie sich an einen Fachanwalt f\"ur Sozialrecht oder die Unabh\"angige Patientenberatung (UPD, 0800 011 77 22).}}
\end{center}


\end{document}
